We start our discussion from Lemma~\eqref{lem:backward-newtons-formula-for-powers},
which yields backward Newton's interpolation formula for powers.
For example,
\begin{align*}
  n^3 &= 0 \tbinom{n-0}{0}  + 1 \tbinom{n-0}{1}  + 6 \tbinom{n-0}{2}  + 6 \tbinom{n-0}{3}, \\
  n^3 &= 1 \tbinom{n-1}{0}  + 7 \tbinom{n-1}{1}  + 12 \tbinom{n-1}{2} + 6 \tbinom{n-1}{3}, \\
  n^3 &= 8 \tbinom{n-2}{0}  + 19 \tbinom{n-2}{1} + 18 \tbinom{n-2}{2} + 6 \tbinom{n-2}{3}, \\
  n^3 &= 27 \tbinom{n-3}{0} + 37 \tbinom{n-3}{1} + 24 \tbinom{n-3}{2} + 6 \tbinom{n-3}{3}.
\end{align*}

Thus, the transition to formula for sums of powers is quite straightforward,
\begin{align*}
  \KnuthRFoldSum{1}{n}{m}
  = \sum_{k=0}^{m} \left[ \sum_{j=1}^{n} \binom{j-t+k-1}{k} \right] \nabla^{k} t^m.
\end{align*}
Therefore, by setting $r=0$ into Lemma~\eqref{lem:backward-hockey-stick-identity}
yields closed form of ordinary sums of powers
\begin{proposition}[Ordinary sums of powers]
  \label{prop:ordinary-sums-of-powers}
  For non-negative integers $n,m$, and an arbitrary integer $t$
  \begin{align*}
    \KnuthRFoldSum{1}{n}{m}
    = \sum_{k=0}^{m}
    \left[
      \binom{n-t+k}{k+1} - \binom{k-t}{k+1}
    \right] \nabla^{k} t^m.
  \end{align*}
\end{proposition}
% We may see that the basis $\binom{j-t+r}{k+r}$ for repeatedly
% applied Lemma~\eqref{lem:forward-hockey-stick-identity}
% was indeed chosen correctly, because we derived closed form for sums of powers
% with ease, simply setting parameter for $r$.
% For example,
% \input{examples/01_ordinary_sums_of_cubes.tex}
% From example above, we may observe binomial coefficients of negative upper
% index, which may looking unfamiliar.
% However, there is no reason to be afraid of, because negative
% upper index binomial coefficients are completely deterministic,
% because $\binom{-n}{k} = (-1)^{k} \binom{k+n-1}{k}$.
% By setting $t=0$ into
% Proposition~\eqref{prop:ordinary-sums-of-powers}
% yields one more well-known identity for sums of powers,
% mentioned in~\cite[p. 190]{graham1994concrete}, and in~\cite{pfaff2007deriving}
% \begin{corollary}
%   For non-negative integers $n,m$
%   \begin{align*}
%     \KnuthRFoldSum{1}{n}{m} = \sum_{j=0}^{m} \binom{n+1}{j+1} \Delta^{j} 0^m.
%   \end{align*}
% \end{corollary}
% For example,
% \begin{example} For non-negative integer $n$, we have
  \begin{align*}
    \KnuthRFoldSum{1}{n}{1} &= 1 \tbinom{n+1}{1},
    \\
    \KnuthRFoldSum{1}{n}{1} &= 0 \tbinom{n+1}{1} + 1 \tbinom{n+1}{2},
    \\
    \KnuthRFoldSum{1}{n}{2} &= 0 \tbinom{n+1}{1} + 1 \tbinom{n+1}{2}
    + 2  \tbinom{n+1}{3},
    \\
    \KnuthRFoldSum{1}{n}{3} &= 0 \tbinom{n+1}{1} + 1 \tbinom{n+1}{2}
    + 6  \tbinom{n+1}{3} + 6   \tbinom{n+1}{4},
    \\
    \KnuthRFoldSum{1}{n}{4} &= 0 \tbinom{n+1}{1} + 1 \tbinom{n+1}{2}
    + 14 \tbinom{n+1}{3} + 36  \tbinom{n+1}{4} + 24  \tbinom{n+1}{5},
    \\
    \KnuthRFoldSum{1}{n}{5} &= 0 \tbinom{n+1}{1} + 1 \tbinom{n+1}{2}
    + 30 \tbinom{n+1}{3} + 150 \tbinom{n+1}{4} + 240 \tbinom{n+1}{5}
    + 120 \tbinom{n+1}{6}.
  \end{align*}
\end{example}

% From the example above, we observe
% the coefficients $0,1,2,0, 1, 6, 6, 0, 1, 14, 36, 24, \dots$,
% are defined by Stirling numbers of the second kind $a_k = k! \stirlingii{n}{k}$,
% which is the sequence
% \href{https://oeis.org/A131689}{\texttt{A131689}}
% in the OEIS~\cite{sloane2003line}.
% This leads us to the notion, that
% Stirling numbers of the second kind $\stirlingii{n}{k}$
% recover naturally from forward Newton's interpolation formula
% for powers.
% By setting $t=1$ into
% Proposition~\eqref{prop:ordinary-sums-of-powers}, we get
% another well-known special case
% \begin{corollary}
%   For non-negative integers $n,m$,
%   \begin{align*}
%     \KnuthRFoldSum{1}{n}{m} = \sum_{j=0}^{m} \binom{n}{j+1} \Delta^{j} 1^m,
%   \end{align*}
% \end{corollary}
% mentioned, for instance, in~\cite{cereceda2022sums}.
% For example,
% \begin{example} For non-negative integer $n$, we have
  \begin{align*}
    \KnuthRFoldSum{1}{n}{0} &= 1 \tbinom{n}{1},
    \\
    \KnuthRFoldSum{1}{n}{1} &= 1 \tbinom{n}{1} + 1 \tbinom{n}{2},
    \\
    \KnuthRFoldSum{1}{n}{2} &= 1 \tbinom{n}{1} + 3 \tbinom{n}{2}
    +2   \tbinom{n}{3},
    \\
    \KnuthRFoldSum{1}{n}{3} &= 1 \tbinom{n}{1} + 7 \tbinom{n}{2}
    +12  \tbinom{n}{3} + 6  \tbinom{n}{4},
    \\
    \KnuthRFoldSum{1}{n}{4} &= 1 \tbinom{n}{1} +15 \tbinom{n}{2}
    +50  \tbinom{n}{3} +60  \tbinom{n}{4} + 24 \tbinom{n}{5},
    \\
    \KnuthRFoldSum{1}{n}{5} &= 1 \tbinom{n}{1} +31 \tbinom{n}{2}
    +180 \tbinom{n}{3} +390 \tbinom{n}{4} + 360 \tbinom{n}{5}
    + 120 \tbinom{n}{6}.
  \end{align*}
\end{example}

% The coefficients $1,1,1,1,3,2,1,7,12,6,1,15, \dots$ is the sequence
% \href{https://oeis.org/A028246}{\texttt{A028246}}
% in the OEIS~\cite{sloane2003line}.
% It is interesting to note that the work~\cite{cereceda2022sums}
% provide formulas for sums of powers
% in terms of generalized Stirling numbers of
% the second kind $\stirlingii{k}{j}_{r}$,
% that is
% \begin{proposition}[Cereceda (2022)]
%   For non-negative integers $n,k$
%   \begin{align}
%     \label{eq:cereceda-sums-1}
%     \KnuthRFoldSum{1}{n}{k} &= \sum_{j=0}^{k} j!
%     \left[
%       \binom{n+1-r}{j+1} + (-1)^j \binom{r+j-1}{j+1}
%     \right] \stirlingii{k}{j}_{r},
%   \end{align}
%   And
%   \begin{align}
%     \label{eq:cereceda-sums-2}
%     \KnuthRFoldSum{1}{n}{k} &= \sum_{j=0}^{k} j!
%     \left[
%       \binom{n+1+r}{j+1} + \binom{r+1}{j+1}
%     \right] \stirlingii{k}{j}_{-r}.
%   \end{align}
% \end{proposition}
% Formula~\eqref{eq:cereceda-sums-1}
% is a special case of
% Proposition~\eqref{prop:ordinary-sums-of-powers}
% with setting $t \rightarrow r$, and upper negated binomial coefficient
% $\binom{1-t}{k+1}$, by using upper negation formula
% from Lemma~\eqref{lem:upper-binomial-negation}.
% We may observe that forward finite difference of powers $\Delta^{k} t^m $ can be expressed
% in terms of generalized Stirling numbers of the second kind, that is
% \begin{align}
%   \label{eq:gen-stirling-finit-difference}
%   \Delta^{k} t^m = k! \stirlingii{m}{k}_{t}.
% \end{align}
% This leads us to the observation, that
% generalized Stirling numbers of the second kind $\stirlingii{n}{k}_t$
% recover naturally from forward Newton's interpolation formula
% for powers, evaluated in arbitrary integer $t$.
% Now, reviewing Formula~\eqref{eq:cereceda-sums-2}, we notice that it is a special case
% of Proposition~\eqref{prop:ordinary-sums-of-powers}
% with setting $t \rightarrow -r$, and by
% using Equation~\eqref{eq:gen-stirling-finit-difference}.
% For example, by setting $t=4$ into
% Proposition~\eqref{prop:ordinary-sums-of-powers}
% yields rather unusual formulas for sums of powers
% \begin{example} For non-negative integer $n$, we have
  \begin{align*}
    \KnuthRFoldSum{1}{n}{0} &= 1  \left( \tbinom{n-3}{1}
    + \tbinom{3}{1}  \right), \\
    \KnuthRFoldSum{1}{n}{1} &= 4  \left( \tbinom{n-3}{1}
    + \tbinom{3}{1}  \right)  + 1  \left( \tbinom{n-3}{2}
    - \tbinom{4}{2}  \right), \\
    \KnuthRFoldSum{1}{n}{2} &= 16 \left( \tbinom{n-3}{1}
    + \tbinom{3}{1}  \right)  + 9  \left( \tbinom{n-3}{2}
    - \tbinom{4}{2}  \right)
    + 2  \left( \tbinom{n-2}{3} + \tbinom{5}{3}  \right), \\
    \KnuthRFoldSum{1}{n}{3} &= 64 \left( \tbinom{n-3}{1}
    + \tbinom{3}{1}  \right)  + 61 \left( \tbinom{n-3}{2}
    - \tbinom{4}{2}  \right)
    + 30 \left( \tbinom{n-3}{3} + \tbinom{5}{3}  \right)
    + 6 \left( \tbinom{n-3}{4} - \tbinom{6}{4}  \right).
  \end{align*}
\end{example}

% The coefficients $1,4,1,16,9,\dots$ is the sequence
% \href{https://oeis.org/A391633}{\texttt{A391633}}
% in the OEIS~\cite{sloane2003line}.
% As we have all required identities in our disposal,
% consider the case of double sums of powers.
% Surely, we derive its formula with ease, and gracefully,
% by setting $r=1$ into hockey-stick identity from
% Lemma~\eqref{lem:forward-hockey-stick-identity},
% with additional correction terms $\KnuthRFoldSum{r}{n}{0}$.
% Therefore,
% \begin{proposition}[Double sums of powers]
%   \label{prop:double-sums-of-powers}
%   For non-negative integers $n,m$, and an arbitrary integer $t$
%   \begin{align*}
%     \KnuthRFoldSum{2}{n}{m}
%     = \sum_{k=0}^{m}
%     \left[
%       \binom{n-t+2}{k+2}
%       - \binom{1-t}{k+1} \KnuthRFoldSum{1}{n}{0}
%       - \binom{2-t}{k+2} \KnuthRFoldSum{0}{n}{0}
%     \right] \Delta^{k} t^m.
%   \end{align*}
% \end{proposition}
% Note that $\KnuthRFoldSum{1}{n}{0}=n$, and $\KnuthRFoldSum{0}{n}{0}=1$.
% By applying upper binomial negation formula
% from Lemma~\eqref{lem:upper-binomial-negation}
% yields
% \begin{proposition}[Double sums of powers negated]
%   \label{prop:double-sums-of-powers-negated}
%   For non-negative integers $n,m$, and an arbitrary integer $t$
%   \begin{align*}
%     \KnuthRFoldSum{2}{n}{m}
%     = \sum_{k=0}^{m}
%     \left[
%       \binom{n-t+2}{k+2} + (-1)^k \binom{k+t-1}{k+1} n
%       + (-1)^{k+1} \binom{k+t-1}{k+2}
%     \right] \Delta^{j} t^m.
%   \end{align*}
% \end{proposition}
% For example, by setting $t=5$ into
% Proposition~\eqref{prop:double-sums-of-powers-negated}
% yields
% \begin{example} For non-negative integer $n$, we have
  \begin{align*}
    \KnuthRFoldSum{2}{n}{0}
    &= 1 \left( \tbinom{n-3}{2} + \tbinom{4}{1} n - \tbinom{4}{2} \right),
    \\
    \KnuthRFoldSum{2}{n}{1}
    &= 5 \left( \tbinom{n-3}{2} + \tbinom{4}{1} n - \tbinom{4}{2} \right)
    + 1 \left( \tbinom{n-3}{3} - \tbinom{5}{2} n + \tbinom{5}{3} \right),
    \\
    \KnuthRFoldSum{2}{n}{2}
    &= 25 \left( \tbinom{n-3}{2} + \tbinom{4}{1} n - \tbinom{4}{2} \right)
    + 11 \left( \tbinom{n-3}{3} - \tbinom{5}{2} n + \tbinom{5}{3} \right)
    + 2 \left( \tbinom{n-3}{4} + \tbinom{6}{3} n - \tbinom{6}{4} \right),
    \\
    \KnuthRFoldSum{2}{n}{3}
    &=
    125 \left( \tbinom{n-3}{2} + \tbinom{4}{1} n - \tbinom{4}{2} \right)
    + 91 \left( \tbinom{n-3}{3} - \tbinom{5}{2} n + \tbinom{5}{3} \right)
    + 36 \left( \tbinom{n-3}{4} + \tbinom{6}{3} n - \tbinom{6}{4} \right) \\
    &+ 6 \left( \tbinom{n-3}{5} - \tbinom{7}{3} n + \tbinom{7}{4} \right).
  \end{align*}
\end{example}

% The coefficients $1,5,1,25,11,2,\dots$ is the sequence
% \href{https://oeis.org/A391635}{\texttt{A391635}}
% in the OEIS~\cite{sloane2003line}.
% Therefore, generalized formula for $r$-fold
% sums of powers is straightforward,
% repeating the same procedure as we have done for
% $\KnuthRFoldSum{1}{n}{m}$, and
% $\KnuthRFoldSum{2}{n}{m}$
% \begin{theorem}[Multifold sums of powers]
%   \label{thm:multifold-sums-of-powers}
%   For non-negative integers $r,n,m$, and an arbitrary integer $t$
%   \begin{align*}
%     \KnuthRFoldSum{r}{n}{m}
%     = \sum_{k=0}^{m}
%     \left[
%       \binom{n-t+r}{k+r}
%       -
%       \sum_{s=0}^{r-1}
%       \binom{r-s-t}{k+r-s}
%       \KnuthRFoldSum{s}{n}{0}
%     \right] \Delta^{k} t^{m}.
%   \end{align*}
% \end{theorem}
% Which is quite remarkable in its pure binomial form,
% by using Lemma~\eqref{lem:multifold-sum-of-zero-powers}
% \begin{proposition}[Multifold binomial sums of powers]
%   \label{prop:multifold-binomial-sums-of-powers}
%   For non-negative integers $r,n,m$, and an arbitrary integer $t$
%   \begin{align*}
%     \KnuthRFoldSum{r}{n}{m}
%     = \sum_{k=0}^{m}
%     \left[
%       \binom{n-t+r}{k+r}
%       -
%       \sum_{s=0}^{r-1}
%       \binom{r-s-t}{k+r-s}
%       \binom{s+n-1}{s}
%     \right] \Delta^{k} t^{m}.
%   \end{align*}
% \end{proposition}
% Explicitly,
% \begin{align*}
%   \KnuthRFoldSum{r}{n}{m}
%   = \sum_{k=0}^{m}
%   \bigg[
%     \tbinom{n-t+r}{k+r}
%     -
%     \tbinom{r-t}{k+r}
%     \tbinom{n-1}{0}
%     -
%     \tbinom{r-1-t}{k+r-1}
%     \tbinom{n}{1}
%     -
%     \cdots
%     -
%     \tbinom{2-t}{k+2}
%     \tbinom{r+n-3}{r-2}
%     -
%     \tbinom{1-t}{k+1}
%     \tbinom{r+n-2}{r-1}
%   \bigg] \Delta^{k} t^{m}.
% \end{align*}
% For example,
% \input{examples/06_multifold_binomial_sums_of_powers_t_0_3.tex}
% These formulas are polynomials in $n$, with integer coefficients
% which are forward finite differences of powers, evaluated at
% arbitrary integer $t$.
% For example, for integer $0 \leq t \leq 4$, we have
% \begin{example}
  For $r=1$ and $m=3$, Proposition~\ref{prop:multifold-binomial-sums-of-powers}
  yields
  \begin{align*}
    t=0:\quad
    \KnuthRFoldSum{1}{n}{3}
    &=
    n\cdot0
    +\tfrac{1}{2}(n+n^2)\cdot1
    +\tfrac{1}{6}(-n+n^3)\cdot6
    +\tfrac{1}{24}(2n-n^2-2n^3+n^4)\cdot6,
    \\
    t=1:\quad
    \KnuthRFoldSum{1}{n}{3}
    &=
    n\cdot1
    +\tfrac{1}{2}(-n+n^2)\cdot7
    +\tfrac{1}{6}(2n-3n^2+n^3)\cdot12 \\
    &+\tfrac{1}{24}(-6n+11n^2-6n^3+n^4)\cdot6,
    \\
    t=2:\quad
    \KnuthRFoldSum{1}{n}{3}
    &=
    n\cdot8
    +\tfrac{1}{2}(-3n+n^2)\cdot19
    +\tfrac{1}{6}(11n-6n^2+n^3)\cdot18 \\
    &+\tfrac{1}{24}(-50n+35n^2-10n^3+n^4)\cdot6,
    \\
    t=3:\quad
    \KnuthRFoldSum{1}{n}{3}
    &=
    n\cdot27
    +\tfrac{1}{2}(-5n+n^2)\cdot37
    +\tfrac{1}{6}(26n-9n^2+n^3)\cdot24 \\
    &+\tfrac{1}{24}(-154n+71n^2-14n^3+n^4)\cdot6,
    \\
    t=4:\quad
    \KnuthRFoldSum{1}{n}{3}
    &=
    n\cdot64
    +\tfrac{1}{2}(-7n+n^2)\cdot61
    +\tfrac{1}{6}(47n-12n^2+n^3)\cdot30 \\
    &+\tfrac{1}{24}(-342n+119n^2-18n^3+n^4)\cdot6.
  \end{align*}
\end{example}
% \begin{example}
%   \begin{align*}
%     t=0: \quad \KnuthRFoldSum{1}{n}{3}
%     &=
%     \\
%     t=1: \quad \KnuthRFoldSum{1}{n}{3}
%     &=
%     \\
%     t=2: \quad \KnuthRFoldSum{1}{n}{3}
%     &=
%     \\
%     t=3: \quad \KnuthRFoldSum{1}{n}{3}
%     &=
%     \\
%     t=4: \quad \KnuthRFoldSum{1}{n}{3}
%     &=
%   \end{align*}
% \end{example}

% \begin{array}{l}
%   n \\
%   \frac{1}{2} n (n+1) \\
%   \frac{1}{6} (n-1) n (n+1) \\
%   \frac{1}{24} (n-2) (n-1) n (n+1) \\
% \end{array}

% Where the coefficients
% \begin{itemize}
%   \item For $t=0$ are $0, 1, 6, 6$ belong to the sequence \href{https://oeis.org/A131689}{\texttt{A131689}} in the OEIS.
%   \item For $t=1$ are $1, 7, 12, 6$ belong to the sequence \href{https://oeis.org/A028246}{\texttt{A028246}} in the OEIS.
%   \item For $t=2$ are $8, 19, 18, 6$ belong to the sequence \href{https://oeis.org/A038719}{\texttt{A038719}} in the OEIS.
%   \item For $t=3$ are $27, 37, 24, 6$ belong to the sequence \href{https://oeis.org/A391552}{\texttt{A391552}} in the OEIS.
%   \item For $t=4$ are $64, 61, 30, 6$ belong to the sequence \href{https://oeis.org/A391633}{\texttt{A391633}} in the OEIS.
% \end{itemize}
